\documentclass[11pt,a4paper]{article}
\usepackage[utf8]{inputenc}
\usepackage[T1]{fontenc}
\usepackage{amsfonts}
\usepackage[french]{babel}
\frenchbsetup{StandardLists=true}
\usepackage{enumitem}
\usepackage{amssymb}
\usepackage{amsmath}
\usepackage{parskip}
\usepackage[left=1.5cm,right=1.5cm,top=1.5cm,bottom=1.5cm]{geometry}
\usepackage{multirow}
\usepackage[dvipsnames]{xcolor}

\usepackage[T1]{fontenc}
\usepackage{fouriernc}

\usepackage{colortbl}

\usepackage{setspace}
\singlespacing

\usepackage{fancyhdr}
\pagestyle{fancy}
\renewcommand\headrulewidth{1pt}
\fancyhead[L]{Lycée Denis Diderot }
\fancyhead[R]{Classe de 1 Bac Pro }
\setcounter{page}{17}\fancyfoot[C]{page \thepage}
\fancyfoot[R]{Année 2024-2025}
\fancyfoot[L]{Alexis RUIZ}
\usepackage{multirow,tabularx,booktabs}

\usepackage{eurosym}

\renewcommand{\arraystretch}{1.8}
\newcolumntype{C}[1]{>{\centering\arraybackslash }b{#1}}

\frenchbsetup{StandardLists=true}
\setlength{\parskip}{2mm}

\begin{document}

\begin{center}
\LARGE{Barème - Évaluation Puissance et énergie}\\[3mm]
\large{Total : 20 points + 1 point bonus}
\vspace {4 mm}

\hrule 

\vspace {4 mm}

\begin{tabular}{|p{11cm}|C{1.5cm}|C{1.5cm}|}
\hline
\rowcolor{gray!30} \textbf{Exercice 1 : QCM} & \textbf{Détail} & \textbf{Total} \\
\hline
Question 1 : joule ET wattseconde (0,5 pt chacune) & 1 pt & \multirow{3}{*}{\textbf{3 pts}} \\
\cline{1-2}
Question 2 : $P=\dfrac{E}{t}$ & 1 pt & \\
\cline{1-2}
Question 3 : $E= P \times t$ & 1 pt & \\
\hline
\end{tabular}

\vspace{3mm}

\begin{tabular}{|p{11cm}|C{1.5cm}|C{1.5cm}|}
\hline
\rowcolor{gray!30} \textbf{Exercice 2 : Calculer une énergie} & \textbf{Détail} & \textbf{Total} \\
\hline
\textbf{Q1.} Durée 8h (0,5) + Formule (0,5) + AN (0,5) + Résultat 112 Wh (1) & 2,5 pts & \multirow{3}{*}{\textbf{5 pts}} \\
\cline{1-2}
\textbf{Q2.} Conversion en kWh : calcul (1) + résultat 0,112 kWh (0,5) & 1,5 pt & \\
\cline{1-2}
\textbf{Q3.} Coût : calcul (0,5) + résultat 0,02 \euro ~(0,5) & 1 pt & \\
\hline
\end{tabular}

\vspace{3mm}

\begin{tabular}{|p{11cm}|C{1.5cm}|C{1.5cm}|}
\hline
\rowcolor{gray!30} \textbf{Exercice 3 : Voiture électrique} & \textbf{Détail} & \textbf{Total} \\
\hline
\textbf{Q1.} 45 W = puissance nominale (1) + 12 V = tension nominale (1) & 2 pts & \multirow{3}{*}{\textbf{6 pts}} \\
\cline{1-2}
\textbf{Q2.} $2 \times 45$ (0,5) + $4 \times 21$ (0,5) + $10 \times 5$ (0,5) + somme 224 W (0,5) + unité (0,5) & 2,5 pts & \\
\cline{1-2}
\textbf{Q3.} Formule $I = \frac{P}{U}$ (0,5) + AN (0,5) + résultat 18,7 A (0,5) & 1,5 pt & \\
\hline
\end{tabular}

\vspace{3mm}

\begin{tabular}{|p{11cm}|C{1.5cm}|C{1.5cm}|}
\hline
\rowcolor{gray!30} \textbf{Exercice 4 : Tableau (0,25 pt par case)} & \textbf{Détail} & \textbf{Total} \\
\hline
Ligne 1 - Chargeur : E = 5,5 Wh (0,5) + Coût = 0,44 \euro ~(0,5) & 1 pt & \multirow{5}{*}{\textbf{6 pts}} \\
\cline{1-2}
Ligne 2 - Lampe : E = 30 Wh (0,5) + Coût = 2,42 \euro ~(0,5) & 1 pt & \\
\cline{1-2}
Ligne 3 - Ordinateur : P = 114 W (0,5) + E annuelle = 209 kWh (0,5) & 1 pt & \\
\cline{1-2}
Ligne 4 - TV : P = 100 W (0,5) + E jour = 301 Wh (0,5) + Coût = 24,20 \euro ~(0,5) & 1,5 pt & \\
\cline{1-2}
Ligne 5 - Four : t = 0,5 h (0,5) + E annuelle = 365 kWh (0,5) + Coût = 80,30 \euro ~(0,5) & 1,5 pt & \\
\hline
\end{tabular}

\vspace{3mm}

\begin{tabular}{|p{11cm}|C{1.5cm}|C{1.5cm}|}
\hline
\rowcolor{gray!30} \textbf{Énigme Bonus} & \textbf{Détail} & \textbf{Total} \\
\hline
Réponse correcte à l'énigme & +1 pt & \textbf{+1 pt} \\
\hline
\end{tabular}

\vspace{5mm}

\end{center}

\end{document}